\documentclass[11pt]{article}
\usepackage[margin=1in]{geometry}
\usepackage[T1]{fontenc}
\usepackage[utf8]{inputenc}
\usepackage{lmodern}
\usepackage{setspace}
\usepackage{hyperref}
\usepackage{enumitem}
\usepackage{booktabs}
\usepackage{longtable}
\usepackage{array}
\usepackage{xcolor}

\hypersetup{
  colorlinks=true,
  linkcolor=blue,
  citecolor=blue,
  urlcolor=blue,
  pdftitle={SignalLock Research Proposal},
  pdfauthor={Sagar Kishore Kumar}
}

\title{\textbf{SignalLock: A Dual-Layer OSINT-Calibrated Framework for Password Predictability and Context-Aware Enterprise Authentication Hardening}}
\author{Sagar Kishore Kumar \\
Independent Research Project \\
\texttt{47300709+sagarkishore-7@users.noreply.github.com}}
\date{\today}

\begin{document}
\maketitle
\onehalfspacing

\begin{abstract}
Password-based authentication remains a major source of organizational risk, yet most password strength estimation mechanisms remain context-free. They are typically designed to score a password without modeling how a user's publicly visible digital footprint may change the probability that the password can be guessed in a targeted attack. At the same time, recent password research has shown that personal information, prior password reuse, and auxiliary attributes inferred from public identifiers can materially improve targeted password guessing, while defensive systems remain comparatively underdeveloped. This proposal introduces \emph{SignalLock}, a defensive, privacy-conscious framework that separates two related but distinct variables: \emph{public exposure risk}, which captures how targetable an account becomes due to organization-approved or consented OSINT, and \emph{password predictability risk}, which estimates how guessable a candidate password becomes when conditioned on that exposure. The system then combines both signals into calibrated authentication hardening recommendations such as warning, rejection, step-up authentication, or mandatory multi-factor enforcement.

The proposed research makes three main contributions. First, it introduces a dual-layer threat model that explicitly distinguishes user exposure from password weakness, addressing a common conceptual flaw in several emerging OSINT-aware defensive ideas. Second, it proposes a practical risk-scoring architecture that can be evaluated on synthetic, anonymized, or otherwise ethical datasets without relying on exploit-ready offensive tooling. Third, it adds an explanation layer designed to help security teams and end users understand why specific public features and password patterns increase risk. The project is aimed at producing a master's thesis, a conference-ready paper, and a working prototype suitable for controlled enterprise evaluation.
\end{abstract}

\section{Introduction}
Passwords continue to serve as a dominant authentication factor across enterprise and consumer systems despite long-standing evidence of their usability and security limitations. Studies on password policies, reuse, and password strength meters have repeatedly shown that users adapt to constraints in predictable ways, often creating passwords that are memorable but structurally weak or vulnerable to reuse-based attacks~\cite{egelman2011passwords,habib2017blacklists,nistsp80063b}. Modern password modeling research has advanced from Markov and PCFG approaches to neural, GAN-based, and LLM-based models that more effectively capture password distributions and targeted reuse behaviors~\cite{weir2009pcfg,melicher2016neural,wheeler2016zxcvbn,passgpt2024,ptrguess2024}.

Recent targeted password guessing work has also demonstrated that auxiliary user information such as personally identifiable information (PII), email-derived attributes, and previously leaked passwords can substantially improve targeted guessing success~\cite{targuess2016,universalmachines2024,dtu2026email}. However, the defensive ecosystem has not kept pace. Most deployed password strength meters remain context-free, and even newer targeted password strength proposals focus narrowly on PII-aware scoring rather than a broader, privacy-conscious authentication hardening workflow~\cite{piipsm2023,wang2023psm}. In practice, enterprise defenders increasingly care about a larger question: which accounts are both publicly exposed and password-vulnerable enough to merit stronger controls?

This proposal argues that the field needs a defensively oriented framework that can jointly model public exposure and candidate-password risk while keeping them analytically distinct. Without that distinction, systems risk conflating two variables that are correlated but not identical. A highly visible executive may be more attractive to attackers without necessarily choosing a weak password, while an obscure employee may still choose a predictable credential. SignalLock addresses this by introducing a dual-layer model that produces three outputs: an exposure score, a candidate-password risk score, and a calibrated hardening recommendation.

\section{Problem Statement}
Existing password security tools exhibit three important gaps.

\begin{enumerate}[leftmargin=*]
  \item \textbf{Context-free password scoring.} Generic password meters such as rule-based or distribution-based estimators do not account for user-specific public context, even though targeted attackers do.
  \item \textbf{Conflation of exposure and weakness.} Emerging ideas about OSINT-aware password defense often collapse public targetability and password predictability into one score, which weakens both interpretability and scientific validity.
  \item \textbf{Limited defensive deployment models.} Existing targeted guessing research is mainly offensive or measurement-focused; it rarely results in privacy-preserving tools for enterprise hardening, candidate-password guidance, or adaptive authentication policy.
\end{enumerate}

The core research problem is therefore: \emph{How can a defensive system use ethical, organization-approved OSINT to estimate targeted password risk and support context-aware enterprise authentication hardening without creating exploit-ready offensive capability?}

\section{Research Questions}
\begin{enumerate}[leftmargin=*]
  \item \textbf{RQ1:} To what extent does organization-approved or consented public OSINT improve defensive prediction of targeted password risk beyond context-free password strength meters?
  \item \textbf{RQ2:} Does separating exposure risk from password predictability produce better calibrated and more actionable security decisions than a single blended score?
  \item \textbf{RQ3:} Which categories of public attributes (professional, temporal, naming, relational, and platform-presence features) contribute most to exposure risk and to conditional password risk?
  \item \textbf{RQ4:} Can dual-layer explanations provide more actionable feedback to security analysts and users than generic password-strength feedback?
\end{enumerate}

\section{Hypotheses}
\begin{enumerate}[leftmargin=*]
  \item \textbf{H1:} Exposure-aware candidate-password scoring will outperform context-free password meters in identifying high-risk passwords under targeted online-guess budgets.
  \item \textbf{H2:} Explicitly separating public exposure from password predictability will improve calibration and reduce false positives compared to a one-score design.
  \item \textbf{H3:} Professional and temporal public attributes, including public name variants, organizational role signals, and year-like identifiers, will contribute substantially to exposure and conditional password-risk estimates.
  \item \textbf{H4:} Dual-layer explanations will be rated as more actionable by security professionals than generic strength-meter feedback.
\end{enumerate}

\section{Related Work}
\subsection{Password Modeling and Strength Estimation}
Classic password modeling approaches such as PCFGs and Markov models established that human-chosen passwords occupy highly non-uniform distributions~\cite{weir2009pcfg}. Neural approaches later improved modeling accuracy while making client-side strength estimation more practical~\cite{melicher2016neural}. Systems such as zxcvbn demonstrated that lightweight, deployable strength estimation can meaningfully improve user feedback for online-attack scenarios~\cite{wheeler2016zxcvbn}. More recent work has explored large language models and adaptive modeling techniques including PassGPT and Universal Neural-Cracking-Machines, indicating that auxiliary information can help tailor password models to target populations~\cite{passgpt2024,universalmachines2024}.

\subsection{Targeted Password Guessing}
TarGuess introduced a structured framework for targeted online password guessing using PII and sister passwords~\cite{targuess2016}. Later work, including Pass2Edit and PointerGuess, improved the modeling of password reuse and tweaking behaviors~\cite{pass2edit2023,ptrguess2024}. By 2026, research had already begun using richer email-derived user attributes and combined PII-plus-old-password modeling with LLMs~\cite{dtu2026email,passglm2026}. These papers are important because they clarify that public and quasi-public attributes can materially affect targeted guessability, but they mainly optimize offensive effectiveness rather than defensive mitigation.

\subsection{Password Behavior, Breach Alerting, and Measurement}
Password policy research has shown that user responses to complexity rules and blacklists are often predictable and sometimes counterproductive~\cite{egelman2011passwords,habib2017blacklists}. Large-scale measurement studies have documented common password reuse and modification practices across services~\cite{domino2018}. Gossamer demonstrated that live password-related measurement can be performed more safely in real authentication systems, while breach alerting systems from Google and follow-up work such as Might I Get Pwned showed the operational value of targeted defensive feedback around compromised credentials~\cite{gossamer2022,google2019breachalert,migp2022}. SignalLock builds on this defensive tradition rather than the offensive branch of targeted guessing.

\subsection{Gap Summary}
The literature suggests four important gaps:

\begin{itemize}[leftmargin=*]
  \item few systems estimate \emph{targeted password risk} in a defensive, candidate-password-aware way;
  \item few systems clearly separate \emph{public exposure} from \emph{password weakness};
  \item few systems offer explainable, calibrated outputs suitable for enterprise authentication policy;
  \item few systems are designed from the outset around ethical, synthetic, or consent-based evaluation.
\end{itemize}

\section{Proposed Contribution}
SignalLock proposes a dual-layer defensive framework with three core contributions:

\begin{enumerate}[leftmargin=*]
  \item \textbf{Dual-layer risk decomposition.} The framework separately models exposure risk and candidate-password predictability risk, then combines them only at the policy stage.
  \item \textbf{Privacy-conscious targeted scoring.} Instead of generating guesses, the system estimates conditional password vulnerability and maps it to bounded risk classes.
  \item \textbf{Actionable explainability.} The framework produces analyst-facing and user-facing explanations that identify which public attributes and password features drive risk.
\end{enumerate}

\section{Methodology}
\subsection{Phase 1: Threat Model and Feature Schema}
The project will begin by formalizing a defensive threat model for targeted online password attacks. The attacker is assumed to have access to publicly available or previously exposed auxiliary data, but the research will focus only on defensive estimation of resulting risk. A structured feature taxonomy will then be defined, separating:

\begin{itemize}[leftmargin=*]
  \item \textbf{Exposure features:} public role prominence, public platform count, naming richness, public date markers, organization-specific identifier visibility, and related signals.
  \item \textbf{Password-conditioned features:} candidate-password overlap with public attributes, transformability, year-like suffixes, common structure classes, and targeted online-risk approximations.
\end{itemize}

\subsection{Phase 2: Ethical Dataset Construction}
Evaluation will rely primarily on synthetic or otherwise ethical datasets. A synthetic organizational profile generator will be built using realistic role, title, organization, and public-bio templates. Publicly available aggregate corpora, such as name-frequency sources and password breach pattern distributions, will be used only for statistical modeling. No real individuals will be profiled without authorization, and no real passwords will be stored in raw form for the interactive mode.

The data strategy includes:

\begin{itemize}[leftmargin=*]
  \item synthetic employee profiles with plausible public information patterns;
  \item synthetic or consent-based candidate-password datasets for controlled evaluation;
  \item aggregate breach-derived pattern statistics from public research corpora;
  \item optional expert review datasets for explanation actionability studies.
\end{itemize}

\subsection{Phase 3: Exposure Engine}
The exposure engine will ingest organization-approved public profile text or profile snapshots and normalize them into structured attribute vectors. Named-entity extraction, alias detection, date recognition, platform-presence detection, and organization-specific identifier recognition will be used to compute an exposure score. Crucially, this score will not claim to represent password weakness directly. Instead, it will represent public targetability and attribute richness relevant to targeted attacks.

\subsection{Phase 4: Candidate-Password Risk Engine}
The password risk engine will accept a candidate password and estimate its targeted online risk conditioned on the normalized attribute vector. It will combine:

\begin{itemize}[leftmargin=*]
  \item generic strength features,
  \item overlap with public attributes,
  \item likely transformations and suffix patterns,
  \item and conditional pattern probabilities informed by aggregate breach statistics.
\end{itemize}

The system will compare several modeling approaches, such as a gradient-boosted tree classifier and a lightweight calibrated neural baseline, to produce bounded vulnerability classes rather than actual guesses.

\subsection{Phase 5: Policy Engine and Explainability}
The policy engine will combine the exposure score and password risk score into authentication hardening actions. Example actions include: accept, warn, reject, require stronger password, enforce MFA, or trigger step-up authentication. An explanation layer, likely using SHAP for model attributions plus template-based natural language generation, will separately explain:

\begin{itemize}[leftmargin=*]
  \item why the account is publicly exposed,
  \item why the candidate password is risky under that exposure,
  \item and why the final policy action was selected.
\end{itemize}

\subsection{Phase 6: Evaluation}
Evaluation will be conducted at three levels:

\begin{enumerate}[leftmargin=*]
  \item \textbf{Model accuracy and calibration.} Precision, recall, F1, AUROC, and calibration error will be measured for high-risk classification.
  \item \textbf{Comparative benchmarking.} The candidate-password engine will be compared against context-free baselines such as zxcvbn-style scoring.
  \item \textbf{Explanation actionability.} Security professionals will assess whether dual-layer explanations are more actionable than generic strength feedback.
\end{enumerate}

\section{System Architecture}
SignalLock will support two initial operating modes:

\subsection{Audit Mode}
For authorized enterprise defenders. Inputs include an organization-approved roster and public profile data. Outputs include exposure heatmaps, prioritized hardening recommendations, and explainable risk summaries.

\subsection{Interactive Mode}
For password creation or password change. Inputs include a candidate password and a local or organization-approved attribute vector. Outputs include a candidate-password risk score, explanation, and hardening recommendation.

The architecture will consist of:

\begin{itemize}[leftmargin=*]
  \item an ingestion and normalization module,
  \item an exposure engine,
  \item a candidate-password risk engine,
  \item a policy engine,
  \item and an explainability layer.
\end{itemize}

\section{Ethical Considerations}
The project explicitly excludes offensive targeting, exploit-ready password generation, and profiling of unauthorized real individuals. Ethical design choices include:

\begin{itemize}[leftmargin=*]
  \item synthetic-first dataset construction,
  \item use of aggregate breach patterns rather than re-identifiable credential pairs,
  \item organization-approved or consent-based public data only,
  \item no storage of real plaintext passwords in the intended interactive workflow,
  \item and clear access controls and audit logging for enterprise evaluation.
\end{itemize}

This is especially important because the research area sits close to targeted password guessing literature. SignalLock aims to improve defense, not to democratize targeted offensive capability.

\section{Expected Results}
The project is expected to show that:

\begin{enumerate}[leftmargin=*]
  \item exposure-aware candidate-password scoring improves identification of high-risk passwords under targeted online-guess assumptions;
  \item dual-layer modeling improves calibration and interpretability compared to single-score OSINT-aware designs;
  \item certain public features, especially name variants, role visibility, and year-like public identifiers, contribute materially to conditional risk;
  \item analyst-facing explanations based on dual-layer attributions are more actionable than generic strength-meter messages.
\end{enumerate}

\section{Expected Deliverables}
The planned deliverables are:

\begin{itemize}[leftmargin=*]
  \item a thesis-quality research proposal and literature synthesis,
  \item a reproducible synthetic evaluation pipeline,
  \item a prototype implementation supporting audit and interactive modes,
  \item a benchmark report comparing SignalLock against context-free baselines,
  \item and a conference-style paper draft suitable for venues such as SOUPS, ACSAC, RAID, or \emph{Computers \& Security}.
\end{itemize}

\section{Timeline}
\begin{longtable}{>{\raggedright\arraybackslash}p{0.18\textwidth} >{\raggedright\arraybackslash}p{0.74\textwidth}}
\toprule
\textbf{Phase} & \textbf{Planned Work} \\
\midrule
Weeks 1--2 & Finalize threat model, terminology, and feature taxonomy. \\
Weeks 3--5 & Build synthetic profile and candidate-password evaluation datasets. \\
Weeks 6--7 & Implement exposure ingestion and normalization pipeline. \\
Weeks 8--10 & Implement candidate-password risk engine and baseline comparisons. \\
Weeks 11--12 & Add calibration, policy engine, and explanation layer. \\
Weeks 13--15 & Build CLI and initial dashboard or API workflow. \\
Weeks 16--18 & Run evaluation, ablation studies, and expert actionability review. \\
Weeks 19--20 & Prepare thesis chapter or conference-style paper draft. \\
\bottomrule
\end{longtable}

\section{Conclusion}
SignalLock proposes a novel and defensible research direction at the intersection of password security, OSINT, and adaptive authentication. Its main conceptual contribution is to separate public exposure from password predictability while still allowing both to inform authentication hardening. This resolves an important methodological weakness in many intuitive OSINT-aware password-defense ideas and creates a clearer path to implementation, evaluation, and publication. Because the project is designed around ethical data use, bounded defensive outputs, and explainable decisions, it is well suited for a master's thesis, a prototype system, and a conference paper.

\bibliographystyle{plain}
\bibliography{references}

\end{document}
